
%% --------------------------------------------------------
\chapter{Хаотична інфляція Лінде}
\label{sec:chaotic}
%% --------------------------------------------------------

Лінде~\cite{linde_1983} запропонував радикально простий погляд
на початкові умови: значення $\phi$ у різних точках раннього
Всесвіту є хаотичним~--- немає потреби в особливому «стартовому»
стані. Найпростіший потенціал:
\begin{equation*}
  V(\phi) = \tfrac{1}{2}m^2\phi^2.
\end{equation*}

Де $\phi$ велике і $\varepsilon \ll 1$~--- там інфляція іде довго
і виникає величезна причинно зв'язана «бульбашка». Де $\phi$ мале~---
інфляція швидко завершується. Наш спостережуваний Всесвіт~---
одна з таких бульбашок.

\textbf{Структура всередині однієї бульбашки.} Коли інфляція
в нашій бульбашці завершилась, простір всередині неї вже був
практично однорідним~--- розширення загладило всі початкові
неоднорідності $\phi$. Але «практично» не означає «абсолютно»:
залишились крихітні квантові флуктуації, заморожені під час
інфляції. Саме вони стали насінням галактик і є тим, що ми
бачимо на карті CMB.

Всередині однієї бульбашки існує єдиний причинно зв'язаний
простір-час~--- один Всесвіт. Але його розмір після $N \sim 60$
$e$-folds становить $\sim e^{60}$ разів більше за наш
спостережуваний горизонт ($\sim 46$~Гпк). Тому навіть у
межах «нашої бульбашки» існують області, з якими ми ніколи
не матимемо причинного контакту~--- вони назавжди за
нашим горизонтом подій.

\textbf{Між бульбашками.} Простір між різними бульбашками
продовжує розширюватись інфляційно~--- швидше за світло.
Дві бульбашки, що утворились незалежно, розділені областю
вічної інфляції і принципово не можуть обмінятись сигналом.
Це не просто «дуже далеко»~--- між ними немає спільного
майбутнього світлового конуса. У строгому сенсі загальної теорії відносності вони
належать різним причинно відокремленим регіонам
одного глобального простору-часу, а не різним «всесвітам»~---
хоча в популярній літературі їх так і називають.

\textbf{Вічна інфляція.} Квантові флуктуації інфлатона на
суперхабблових масштабах можуть локально \emph{збільшити} $\phi$,
повернувши його назад на схил, якщо флуктуація достатньо велика.
Класичний зсув поля за один $e$-fold: $|\delta\phi_{\rm class}|
\sim |\dot\phi|/H \approx \sqrt{2\varepsilon}\,M_{\rm Pl}$.
Квантова флуктуація за той самий час: $\delta\phi_{\rm quant}
\sim H/(2\pi)$. Умова «самовідтворення»~---
квантова флуктуація перевищує класичне зміщення:
\begin{equation*}
  \frac{H}{2\pi} \gtrsim \sqrt{2\varepsilon}\,M_{\rm Pl}
  \quad\Leftrightarrow\quad
  H^2 \gtrsim M_{\rm Pl}^2
  \quad\Leftrightarrow\quad
  V(\phi) \gtrsim M_{\rm Pl}^4,
\end{equation*}
де в останньому переході використано $H^2 \approx V/(3M_{\rm Pl}^2)$
і $\varepsilon \lesssim 1$.
У таких областях інфляція ніколи не зупиняється глобально:
поки одна бульбашка завершує інфляцію і термалізується,
сусідні регіони продовжують розширюватись, породжуючи нові
бульбашки. Це \textbf{вічна інфляція}~--- мультивсесвіт
як математичний наслідок, а не постулат.

\textbf{Статус.} Потенціал $\tfrac{1}{2}m^2\phi^2$ передбачає
$r \approx 8/N \approx 0.13$~--- вище поточної межі BICEP/Keck
(табл.~\ref{tab:inflation_models}), тому ця конкретна модель
виключена. Але ідея хаотичних початкових умов і вічної інфляції
залишається живою у більш пологих потенціалах (Старобінський,
хіггсова інфляція).

Разом з тим жодна з «успішних» моделей~--- тих, що проходять
обмеження Planck на $n_s$ і $r$~--- не передбачає нічого
особливого на малих масштабах $k \gtrsim 1\,h/$Мпк. Саме там
первинний спектр залишається практично необмеженим
спостереженнями. І саме там у 2022--2023 роках телескоп JWST
зафіксував аномалії, які стандартний спектр пояснити не може.